\subsection{API設計}%定期実行する処理のAPIもここ
本システムのAPIは，データベース操作および外部サービスとの連携を行うWeb APIとして表\ref{tab:thisAPI}のように設計している．
具体的には，帰宅推奨時刻，提示する時刻，投票情報およびその結果といったデータの登録・取得のみを実行し，画面遷移やユーザへの提示内容はフロントエンド側に委ねる設計とした．
この設計により，帰宅推奨時刻の算出から投票の集計に至るまでの処理を段階的に管理できるように，各段階に対応するAPIを分離している．
なお，フロントエンド側では独自のAPIエンドポイントを介して本APIを呼び出す構成としており，処理の役割分担を明確化している．
また，本APIは外部からの呼び出しを想定したWeb APIとして公開しているため，既存データの上書きは行わず新たなレコードを追加する設計としている．
ただし，ユーザ情報を扱うAPIについては，在室管理システムや投票ツールとの連携を段階的に行う状況を想定し，不足情報を後から補完できる設計としている．


\begin{table}[htb]
  \centering
  \caption{本システムで提供するAPIの分類}
  \label{tab:thisAPI}
  \begin{tabular}{|l|l|}
    \hline
    \textbf{API分類} & \textbf{役割概要}                \\
    \hline
    ユーザ・連携API      & ユーザ情報管理および外部サービス連携           \\
    帰宅推奨時刻管理API    & 帰宅推奨時刻の登録および最新状態の取得          \\
    提示時刻管理API      & 提示する時刻候補と投票締切時刻の登録および最新状態の取得 \\
    投票管理API        & 投票情報の受付および登録                 \\
    投票結果管理API      & 投票集計結果の登録および最新状態の取得          \\
    \hline
  \end{tabular}
\end{table}


ユーザ・連携APIは，ユーザ情報の管理および外部サービスとの連携を目的としている．
本APIでは，ユーザ名による検索は提供しない設計としており，外部サービス側で管理される識別子を用いた検索に限定している．
これにより，フロントエンド側から特定ユーザの在室状況や通知対象を推測できてしまう可能性を低減している．
また，将来的に在室管理システムや投票ツールとの連携を段階的に拡張できるよう，不足情報を後から補完可能な設計としている．


帰宅推奨時刻管理APIは，投票や時刻提示の起点となる帰宅推奨時刻を管理するために用意している．
本APIは定期実行により呼び出され，事前に判定された条件を満たす場合に帰宅推奨時刻を登録する．
提示時刻管理APIや投票処理は，本APIによって記録されたデータを基に動作する構成としており，処理の段階性を明確にしている．
なお，本APIは外部から無制限に呼び出される状況を防ぐため，定期実行を行う GitHub Actions からのみ実行可能な設計としている．
具体的には，API呼び出し時に認証用のキーを要求し，このキーを GitHub Actions の環境変数としてのみ保持し，第三者による不正な実行を防止している．
これにより，帰宅推奨時刻の算出処理が意図しないタイミングで過剰に実行される状況を防ぎ，運用上の安全性を確保している．


提示時刻管理APIおよび投票結果管理APIは，いずれもフロントエンドにおける画面表示に直接関与するAPIである．
本システムでは，これらをまとめて提示状態管理APIとして位置付けている．
提示時刻管理APIは，帰宅推奨時刻が現在時刻に近づいた場合に，画面上で提示する時刻情報を登録する役割を担う．
投票結果管理APIは，現在時刻が投票締切時刻を過ぎた場合に，集計結果を登録する役割を担う．
なお，「現在時刻に近づいた場合」や「投票締切時刻」の判定条件は，利用環境や運用方針に応じて調整可能なパラメータとして実装している．
投票締切後の集計処理については，クライアント側から具体的な投票内容を指定させるのではなく，提示時刻を識別するIDのみを受け取り，集計および結果の保存はAPI内部で完結する設計としている．
この構成により，集計ロジックの一貫性を保つとともに，クライアント実装の簡素化を図っている．
いずれも画面状態の判定に用いるAPIであり，フロントエンドから1分間隔で高頻度に呼び出されるため，最新レコードのみを取得するエンドポイントを設けている．
これにより，フロントエンド側では現在の状態を簡潔に取得でき，不要な履歴データの取得を回避できる．


投票管理APIは，外部サービスからの投票入力を受け付け，投票情報を記録する役割を担う．
さらに，外部サービスとの連携を考慮し，APIではバックエンド内部のユーザ識別子を外部に出さず，外部サービス側で管理される識別子のみで投票処理やDM送信が完結する設計としている．
本システムにおける投票処理は主にSlackのイベントを通じて受信しているが，APIとしても外部サービス固有の識別子を用いて投票情報を受け付ける構成を実装している．
これにより，実装上はSlack連携に依存しつつも，将来的に他の投票手段を導入する際の拡張性を確保している．
